% Generated draft: requires mathematical and source review.
\documentclass[11pt]{article}
\usepackage[margin=25mm]{geometry}
\usepackage{fontspec}
\setmainfont{Latin Modern Roman}
\usepackage{amsmath,amssymb,mathtools,unicode-math}
\setmathfont{Latin Modern Math}
\usepackage{enumitem,needspace,xurl,hyperref}
\setlength{\parindent}{0pt}
\setlength{\parskip}{5pt}
\setlength{\emergencystretch}{3em}
\newcommand{\N}{\mathbb{N}}
\newcommand{\Z}{\mathbb{Z}}
\newcommand{\Q}{\mathbb{Q}}
\newcommand{\R}{\mathbb{R}}
\newcommand{\C}{\mathbb{C}}
\newcommand{\F}{\mathbb{F}}
\newcommand{\A}{\mathbb{A}}
\newcommand{\PP}{\mathbb{P}}
\newcommand{\E}{\mathbb{E}}
\newcommand{\eps}{\varepsilon}
\newcommand{\dd}{\,\mathrm{d}}
\begin{document}
\section*{UnsolvedMath ID: 30005749}
\textbf{Source problem code:} catalogue entry
\subsection{Definitions and mathematical statement}
\textbf{Research question.} How can the underlying proofs be formalised in a suitable abstract system?

Here, ``the underlying proofs'' and ``a suitable abstract system'' retain their unresolved meanings from the supplied source excerpt. The excerpt does not identify the proofs, their assumptions or conclusions, the proposed formal language or deductive rules, or the criteria that make a system suitable. Consequently, it does not determine a precise quantified mathematical statement or definitions of its objects. Specifying these data is necessary before a formalisation problem can be stated precisely.

The accompanying text names the contribution \emph{New applications of proof theory: Greedy algorithms, probability, and proof assistants} and Thomas Powell. This contextual title alone does not identify which proofs are intended or impose mathematical hypotheses. The supplied record describes the question as broad and open; these are source qualifications, not independently verified conclusions.
\subsection{Short English statement}
The source asks how to formalise certain underlying proofs in an appropriate abstract system. It does not specify those proofs or what counts as an appropriate system, so a precise mathematical target cannot be recovered from the supplied excerpt.
\subsection{Sources}
\begin{itemize}
\item Supplied local UnsolvedMath registry record, permanent ID \(30005749\), title \emph{Abstract Formalization of Proof-Theoretic Arguments}, source problem code \texttt{catalogue entry}, in the Oberwolfach Reports Open Problems collection.
\item Supplied report citation and literature-source link labelled ``original OWR report'': \emph{Mathematical Logic: Proof Theory, Constructive Mathematics} (2024), Oberwolfach Reports, pp.\ 3013--3061, DOI \texttt{10.4171/owr/2023/53}. \url{https://doi.org/10.4171/owr/2023/53}. The supplied excerpt associates the question with \emph{New applications of proof theory: Greedy algorithms, probability, and proof assistants}, Thomas Powell; no separate proposer is specified. The registry gives the proposed year as 2023.
\item Literature assessment dated exactly \texttt{2026-08-22}: the supplied curation labels the question open, reports that the workshop report explicitly presents it as a broad open question, and states that no later primary formalisation resolving it was located. This is supplied dated curation, not independently verified here. Its listed checked source is the original OWR report linked above.
\item The registry labels the statement \texttt{corrected\_verified} and reports checking it against the official Oberwolfach report, including preservation of hypotheses, notation, and the open target. This verification claim is supplied metadata and has not been independently checked here.
\end{itemize}
\end{document}
